\subsection{Practical Findings}\label{sec:practical-findings}

\paragraph{Permutation sensitivity.} The per-byte $a_k$ curve depends on the response ordering.
Comparing the $D_{Ca_n}$ ranking of the \permSensNumScenarios{} validation scenarios from Section~\ref{sec:scenario-validation} between \permSensLowPerm{}- and \permSensHighPerm{}-permutation runs: on GPT-2, \permSensGPTMisranked{} of \permSensNumScenarios{} scenarios shift rank between the two settings; on Qwen2.5-3B, \permSensQwenMisranked{} of \permSensNumScenarios{}.
Where scenarios do shift rank, it is between adjacent scenarios with similar $D_{Ca_n}$ values rather than dramatic reorderings.
We use \permSensHighPerm{} permutations throughout for scenario-level experiments; we have not swept intermediate values, so we cannot pinpoint a minimum that is both cheap and reliable.
The coherence term $C$ is permutation-invariant by construction and unaffected.

\paragraph{Boundary handling.} When a tokenizer produces a single token spanning a response and the following separator (e.g., Qwen's \texttt{.\textbackslash n\textbackslash n}), our boundary detector attributes the merged token to the separator; the response's trailing character contributes to its byte count but not to its cross-entropy.
Boundary correctness is verified by \texttt{tests/test\_response\_boundaries.py}.

